\chapter{Практика}


В данной лабораторной работе необходимо произвести расчет дисперсии шума, зная энергию на символ/бит и SNR (задается пользователем в db).
Сгенерировать биты, сформировать символы (0->-1, 1->1), сделать upsampling, добавить АБГШ, на приеме пропустить сигнал через согласованный фильтр,
взять отсчеты в нужные моменты времени (конец каждого символа).

\section{Расчет шума}
$$SNR_{linear} = 10^{SNR/10}$$
$$E_b = \int_{0}^{T}s[i]^2$$

$$\boxed{sigma = \sqrt{\frac{E_b}{SNR_{linear}}}}$$

\section{Формирование сигнала на передатчике}

\begin{figure}[H]
    \centering
    \includegraphics{2.png}
    \caption{Сигнал на передаче}
\end{figure}

\section{Добавление шума}

\begin{figure}[H]
    \centering
    \includegraphics{3.png}
    \caption{Сигнал на приеме}
\end{figure}

\section{Выход согласованного фильтра}

На приеме необходимо понять, какой символ был передан. Эту задачу можно свести к подсчету корреляционных интегралов

$$\int_{0}^{T}z(t)S_1(t)dt >= \int_{0}^{T}z(t)S_2(t)dt$$

Вычислять такие интегралы не совсем удобно, поэтому и был придуман согласованный фильтр.\\

Согласованный фильтр заменяет корреляторы в приемнике, выполняя свертку сигнала с отраженной формой сигнала.
Мы делаем свертку, но она считает то же самое, что и корреляторы. Это позволяет свести подсчеты корреляционного интеграла
к фильтрации линейным фильтром.

$$\boxed{y(t) = \int_{0}^{t}z(t\tau)h(t_0-\tau)d\tau = \alpha\int_{0}^{t}z_1(\tau)S(\tau)d\tau}$$

здесь $h(t_0-\tau)$ - ИХ согласованного фильтра, которая имеет форму передаваемого сигнала, но зеркально отраженного.

\begin{figure}[H]
    \centering
    \includegraphics{4.png}
    \caption{Выход согласованного фильтра}
\end{figure}

Взяв отсчеты в моменты времени окончания символа, РУ может принять решение о передаваемом бите. Если отсчет больше 0, то передавалась единица, а иначе передавался ноль.
Именно в моменты окончания символа выход согласованного фильтра имеет наибольшее отношение сигнал/шум (SNR). Именно такой отсчет нам наибоее точно "скажет", какой символ передавался.
\section{Сравнение теоретического и эмпирического BER}

\begin{figure}[H]
    \centering
    \includegraphics{5.png}
    \caption{Теоретичекий и эмпирический BER}
\end{figure}


\endinput 